% ============================================================================
%  CHƯƠNG 3. PHƯƠNG PHÁP
% ============================================================================
\clearpage
\section{Phương pháp}
\label{sec:method}

\subsection{Kiến trúc hệ thống}

Hình~\ref{fig:arch} là kiến trúc tổng thể của hệ thống, chia thành sáu tầng. Mỗi tầng ứng với một nhóm module
trong gói \code{nodulecxr} và một nhóm script trong \code{code/scripts/}. Dữ liệu đi một chiều từ trên xuống.
Ảnh gốc của ba nguồn được chuẩn hoá thành PNG 896$\times$896 và một file \code{manifest\_*.csv} cho mỗi nguồn.
Mọi mô hình, VLM hay CNN, đều ghi dự đoán theo cùng một định dạng JSONL. Một chương trình đánh giá duy nhất
(\code{make\_report.py}) đọc các file này, chọn ngưỡng trên validation và sinh mọi bảng của báo cáo. Kiểm toán E5
dùng ngưỡng và cấu hình do tầng đánh giá chọn. Nhờ vậy, VLM và CNN được so trên cùng split, cùng ngưỡng và
cùng mã tính chỉ số.

\begin{figure}[H]
\centering
\begin{tikzpicture}[x=1cm,y=1cm, fb/.append style={font=\scriptsize, minimum height=0.8cm},
    bandlbl/.append style={font=\scriptsize}]
  % ---- nền các tầng: tâm tầng i tại y = -1.75 i ----
  \foreach \y in {0,-1.75,-3.5,-5.25,-7.0,-8.75}{\draw[band] (0,\y-0.72) rectangle (15.4,\y+0.72);}
  \node[bandlbl] at (0.1,0)     {\textbf{1. Nguồn dữ liệu}\\Kaggle, NIH, JSRT};
  \node[bandlbl] at (0.1,-1.75) {\textbf{2. Tiền xử lý (E0)}\\\code{data/}, \code{images.py}\\\code{prepare\_*.py}};
  \node[bandlbl] at (0.1,-3.5)  {\textbf{3. Mô hình}\\\code{vlm/}, \code{train/}\\\code{baselines/}};
  \node[bandlbl] at (0.1,-5.25) {\textbf{4. Dự đoán}\\\code{predict\_vlm.py}\\\code{train\_*.py}};
  \node[bandlbl] at (0.1,-7.0)  {\textbf{5. Đánh giá}\\\code{metrics/}, \code{eval/}\\\code{make\_report.py}};
  \node[bandlbl] at (0.1,-8.75) {\textbf{6. Kiểm toán (E5)}\\\code{audit/}\\\code{run\_audit.py}};
  % ---- tầng 1 ----
  \node[fb, text width=3.95cm] (vin)  at (4.9,0)  {\textbf{VinBigData / VinDr-CXR}\\train · val · test (70/10/20)};
  \node[fb, text width=3.55cm] (nih)  at (9.3,0)  {\textbf{NIH ChestX-ray14}\\\code{test\_nih}, ngoài phân phối};
  \node[fb, text width=3.55cm] (jsrt) at (13.5,0) {\textbf{JSRT}\\\code{test\_jsrt}, độ mờ 1--5};
  % ---- tầng 2 ----
  \node[fb, text width=2.6cm] (p1) at (4.225,-1.75)  {DICOM, PNG, raw\\$\rightarrow$ PNG 896$\times$896};
  \node[fb, text width=2.6cm] (p2) at (7.475,-1.75)  {Hợp nhất box\\bằng WBF (IoU 0,4)};
  \node[fb, text width=2.6cm] (p3) at (10.725,-1.75) {\code{manifest\_*.csv}\\nhãn, box, split};
  \node[fb, text width=2.6cm] (p4) at (13.975,-1.75) {\code{data\_stats.py}\\\code{e0\_data.md}, \code{edges}};
  % ---- tầng 3 ----
  \node[fgray,  text width=2.0cm] (c0)  at (3.94,-3.5)  {\textbf{C0} zero-shot\\MedGemma, Qwen};
  \node[fblue,  text width=2.0cm] (c13) at (6.52,-3.5)  {\textbf{C1--C3}\\LoRA MedGemma};
  \node[fblue,  text width=2.0cm] (c4)  at (9.10,-3.5)  {\textbf{C4}\\LoRA Qwen};
  \node[fgreen, text width=2.0cm] (b1)  at (11.68,-3.5) {\textbf{B1}\\DenseNet-121};
  \node[fgreen, text width=2.0cm] (b2)  at (14.26,-3.5) {\textbf{B2}\\YOLOv8};
  % ---- tầng 4 ----
  \node[fb, text width=12.3cm] (pred) at (9.1,-5.25) {\code{outputs/preds/<cfg>.jsonl}, mỗi dòng một ảnh:\\
    \code{image\_id}, \code{split}, \code{config}, \code{seed}, \code{run}, \code{score}, \code{raw\_text}, \code{boxes}};
  % ---- tầng 5 ----
  \node[forange, text width=2.6cm] (e1) at (4.225,-7.0)  {\code{thresholds.json}\\$t_{90}$, $t_{95}$ trên val};
  \node[forange, text width=2.6cm] (e2) at (7.475,-7.0)  {\code{val\_selection.csv}\\chọn cấu hình C4, E5};
  \node[fb,      text width=2.6cm] (e3) at (10.725,-7.0) {\code{summary.csv/.md}\\bảng chính (run 0)};
  \node[fb,      text width=2.6cm] (e4) at (13.975,-7.0) {\code{strata}, \code{comparisons}\\\code{run\_agreement}};
  % ---- tầng 6 ----
  \node[fred, text width=5.7cm] (au) at (5.85,-8.75) {T1 che nốt · T2 che ngẫu nhiên · T3 tráo ảnh\\$\rightarrow$ \code{outputs/audit/*.json} ($\Delta_{\text{occ}}$, $n_{\text{T2}}$)};
  \node[fb, text width=5.7cm] (rep) at (12.35,-8.75) {Bảng và hình trong báo cáo\\(Mục~\ref{sec:data} và Mục~\ref{sec:results})};
  % ---- mũi tên giữa các tầng ----
  \foreach \x in {4.9,9.3,13.5}{\draw[farr] (\x,-0.5) -- (\x,-1.3);}
  \foreach \x in {3.94,6.52,9.10,11.68,14.26}{\draw[farr] (\x,-2.2) -- (\x,-3.05);}
  \foreach \x in {3.94,6.52,9.10,11.68,14.26}{\draw[farr] (\x,-3.95) -- (\x,-4.8);}
  \foreach \x in {4.225,7.475,10.725,13.975}{\draw[farr] (\x,-5.7) -- (\x,-6.55);}
  \foreach \x in {4.225,7.475}{\draw[farr] (\x,-7.45) -- (\x,-8.3);}
  \foreach \x in {10.725,13.975}{\draw[farr] (\x,-7.45) -- (\x,-8.3);}
\end{tikzpicture}
\caption{Kiến trúc hệ thống theo tầng. Cột trái ghi module và script tương ứng trong \code{code/}. Ô xám: VLM
dùng nguyên trạng; ô xanh dương: VLM có LoRA; ô xanh lá: baseline CNN; ô cam: các file chỉ tính trên validation,
dùng để chọn ngưỡng và cấu hình (C4 và cấu hình cho E5 được chọn từ \code{val\_selection.csv}); ô đỏ: kiểm toán E5.}
\label{fig:arch}
\end{figure}

\subsection{Luồng xử lý của trợ lý}

Hình~\ref{fig:pipeline} mô tả luồng xử lý của trợ lý cho một ảnh. Ảnh CXR được chuẩn hoá về 896$\times$896 và
đưa vào VLM cùng một câu hỏi cố định. VLM trả lời bằng JSON gồm kết luận có/không và danh sách box. Song song
với việc sinh văn bản, mô hình được chấm điểm trên hai câu trả lời ép buộc để có điểm liên tục $s(x)$
(Mục~\ref{sec:score}). Parser kiểm tra cú pháp và chuyển câu trả lời thành nhãn và box; điểm $s(x)$ dùng để vẽ
ROC và chọn ngưỡng.

\begin{figure}[H]
\centering
\begin{tikzpicture}[node distance=0.5cm and 0.45cm, fb/.append style={font=\scriptsize}]
  \node[fb, text width=1.6cm, minimum height=1.05cm] (img) {Ảnh CXR\\896$\times$896};
  \node[fgray, right=of img, text width=2.0cm, minimum height=1.05cm] (enc) {Bộ mã hoá ảnh\\(đông cứng)};
  \node[fblue, right=of enc, text width=1.6cm, minimum height=1.05cm] (llm) {LLM\\+ LoRA $r{=}16$};
  \node[fb, right=of llm, text width=2.75cm, minimum height=1.05cm, yshift=0.8cm] (json) {Sinh văn bản\\JSON yes/no + box};
  \node[fb, right=of llm, text width=2.75cm, minimum height=1.05cm, yshift=-0.8cm] (sc) {Chấm điểm ép buộc\\$s(x)\in[0,1]$};
  \node[fb, right=0.4cm of json, text width=1.9cm, minimum height=1.05cm] (parse) {Parser\\nhãn + box};
  \node[forange, right=0.4cm of sc, text width=1.9cm, minimum height=1.05cm] (thr) {So với ngưỡng\\$t_{90}$, $t_{95}$};
  \node[fb, above=0.45cm of enc, text width=4.6cm, xshift=1.2cm] (q) {Câu hỏi cố định: \emph{``Is there a lung nodule or mass\dots?''}};
  \draw[farr] (img) -- (enc); \draw[farr] (enc) -- (llm);
  \draw[farr] (llm.east) -- ++(0.2,0) |- (json.west);
  \draw[farr] (llm.east) -- ++(0.2,0) |- (sc.west);
  \draw[farr] (json) -- (parse); \draw[farr] (sc) -- (thr);
  \draw[farr] (q.south -| llm.north) -- (llm.north);
\end{tikzpicture}
\caption{Luồng xử lý của trợ lý cho một ảnh. Chỉ các ma trận LoRA trong LLM (ô xanh) được huấn luyện; bộ mã
hoá ảnh (ô xám) đông cứng. Nhánh trên cho câu trả lời đọc được và box; nhánh dưới cho điểm liên tục dùng để chọn
ngưỡng trên validation.}
\label{fig:pipeline}
\end{figure}

\subsection{Mô hình}

\paragraph{VLM chính: MedGemma-4B-it.} LoRA $r = 16$, $\alpha = 32$, dropout 0,05, gắn vào các lớp attention
và MLP của LLM; bộ mã hoá ảnh đông cứng. Lý do chọn: (i) là mô hình y khoa mở mới nhất~\cite{sellergren2026medgemma};
(ii) zero-shot nó rơi đúng vào chế độ báo bình thường~\cite{nishio2026zeroshot}, nên câu hỏi ``LoRA có kéo mô
hình ra khỏi chế độ đó không'' có tính thông tin cao; (iii) đầu vào 896$\times$896 trùng với độ phân giải tối ưu
cho CAD nốt~\cite{kim2020matrix}, nên nút thắt nhiều khả năng nằm ở cách gộp đặc trưng và dữ liệu, không ở độ
phân giải. Rủi ro: MedGemma không được huấn luyện mạnh để xuất toạ độ, nên khả năng định vị phải học hoàn toàn
qua LoRA.

\paragraph{VLM thứ hai: Qwen2.5-VL-3B-Instruct.} Dùng cho zero-shot và cho cấu hình LoRA tốt nhất, để so mô
hình y khoa với mô hình đa năng có sẵn khả năng xuất box. Lưu ý: \cite{nishio2026zeroshot} thử Qwen3-VL-8B
(sensitivity 7,8\%), không phải Qwen2.5-VL-3B, nên không so trực tiếp số liệu.

\paragraph{Baseline CNN.}
\begin{itemize}
  \item \textbf{B1 --- classifier:} DenseNet-121~\cite{huang2017densenet} tiền huấn luyện trên CXR
        (thư viện TorchXRayVision~\cite{cohen2022torchxrayvision}), fine-tune nhị phân nốt/khối.
  \item \textbf{B2 --- detector:} YOLOv8~\cite{jocher2023yolov8} huấn luyện trên box nốt/khối, đầu vào 896.
\end{itemize}

\subsection{Các cấu hình}

\begin{table}[H]
\centering
\caption{Các cấu hình được so sánh.}
\label{tab:configs}
\small
\begin{tabularx}{\textwidth}{@{}l L l@{}}
\toprule
\textbf{Mã} & \textbf{Cấu hình} & \textbf{Đầu ra} \\
\midrule
C0 & VLM zero-shot (MedGemma-4B, Qwen2.5-VL-3B) & JSON yes/no + box \\
C1 & LoRA MedGemma-4B, chỉ phân loại & yes/no \\
C2 & LoRA MedGemma-4B, phân loại và định vị & yes/no + box \\
C3 & C2 + cân bằng lớp (oversampling ca dương 1:1, trọng số loss cho token trả lời) & yes/no + box \\
C4 & Cấu hình tốt nhất trong C1--C3, áp lên Qwen2.5-VL-3B & như cấu hình được chọn \\
B1 & DenseNet-121 classifier & xác suất \\
B2 & YOLOv8 detector & box + điểm \\
\bottomrule
\end{tabularx}
\end{table}

\paragraph{Mở rộng, không thuộc cam kết.} C5 --- VLM có hỗ trợ vùng: đưa box đề xuất của B2, hoặc ảnh cắt theo
box, vào prompt để VLM xác nhận và mô tả, nhắm vào giả thuyết gộp toàn cục~\cite{muthyala2026frozen}. C6 --- chia
ô hoặc cắt theo trường phổi để vượt 896\,px hiệu dụng~\cite{sabottke2020resolution,kim2020matrix}.

\subsection{Định dạng huấn luyện và suy luận}

\paragraph{Câu hỏi cố định.}
\begin{itemize}
  \item Định vị (C0, C2--C4): \emph{``Is there a lung nodule or mass on this chest X-ray? If yes, give the
        bounding box(es) in pixel coordinates. Answer in JSON.''}
  \item Chỉ phân loại (C1): \emph{``Is there a lung nodule or mass on this chest X-ray? Answer in JSON.''}
\end{itemize}

\paragraph{Đáp án đích.} Sinh tự động từ nhãn gốc: box của các bác sĩ được hợp nhất bằng
WBF~\cite{solovyev2021wbf}, quy về toạ độ pixel trên ảnh 896$\times$896 dạng $[x_1,y_1,x_2,y_2]$. Ví dụ
ca dương có đáp án \code{\{"nodule": "yes", "boxes": [[412, 230, 468, 281]]\}}, còn ca âm có đáp án
\code{\{"nodule": "no", "boxes": []\}}. Với C1 (chỉ phân loại), đáp án chỉ có khoá \code{nodule}. Không dùng câu
hỏi--đáp do LLM sinh, để tránh nhiễu nhãn. Không đưa ngữ cảnh lâm sàng vào prompt trong điều kiện chính, vì ngữ
cảnh có thể ``lộ'' đáp án~\cite{cusick2026context}.

\paragraph{Điểm liên tục.}\label{sec:score} Để vẽ ROC và chọn ngưỡng, VLM cần một điểm liên tục chứ không chỉ
nhãn. Đồ án so log-xác suất của hai chuỗi trả lời ép buộc $a_{+} = $ \code{\{"nodule": "yes} và $a_{-} = $
\code{\{"nodule": "no} (không có dấu nháy đóng), tính từ token đầu tiên mà hai chuỗi khác nhau:
\begin{equation}
  s(x) = \frac{p_\theta(a_{+} \mid x, q)}{p_\theta(a_{+} \mid x, q) + p_\theta(a_{-} \mid x, q)} = \sigma\big(\log p_\theta(a_{+} \mid x, q) - \log p_\theta(a_{-} \mid x, q)\big).
\end{equation}
Hình~\ref{fig:yesno} minh hoạ cách tính. Ngưỡng $t_{90}$ là ngưỡng nhỏ nhất mà specificity trên validation đạt
$\ge 0{,}90$, tức khoảng 90\% ảnh âm của validation có điểm dưới ngưỡng; $t_{95}$ tính tương tự. Ngưỡng được tính
riêng cho từng bộ (cấu hình, seed, lần chạy), lưu trong \code{thresholds.json} với khoá \code{config|seed|run},
rồi áp nguyên lên cả ba tập test. Điểm $s(x)$ được tính độc lập với văn bản sinh ra, nên vẫn có cả khi câu trả lời
không parse được. CNN dùng điểm gốc của chúng: xác suất sigmoid (B1) và điểm box cao nhất của ảnh (B2).

\begin{figure}[H]
\centering
\begin{tikzpicture}[node distance=0.4cm and 0.5cm]
  \node[fb, text width=2.6cm, minimum height=1.6cm] (ctx) at (1.4,0) {Ảnh $x$ + câu hỏi $q$\\+ tiền tố chung\\\code{\{\dq{}nodule\dq{}: \dq{}}};
  \node[fgreen, text width=2.6cm] (yes) at (4.95,0.75) {tiếp \code{yes}\\$\log p_\theta(a_{+}\mid x,q)$};
  \node[fred, text width=2.6cm] (no) at (4.95,-0.75) {tiếp \code{no}\\$\log p_\theta(a_{-}\mid x,q)$};
  \node[draw=black!60, circle, inner sep=1pt, font=\small] (minus) at (7.0,0) {$-$};
  \node[fb, text width=0.85cm] (sig) at (8.05,0) {$\sigma(\cdot)$};
  \node[fblue, text width=0.8cm] (s) at (9.3,0) {$s(x)$};
  \draw[farr] (ctx.east) -- ++(0.35,0) |- (yes.west);
  \draw[farr] (ctx.east) -- ++(0.35,0) |- (no.west);
  \draw[farr] (yes.east) -| (minus.north);
  \draw[farr] (no.east) -| (minus.south);
  \draw[farr] (minus) -- (sig); \draw[farr] (sig) -- (s);
  % ---- ô phải: chọn ngưỡng trên val ----
  \begin{scope}[xshift=10.55cm, yshift=-1.25cm]
  \begin{axis}[width=5.0cm, height=3.6cm, xmin=0, xmax=1, ymin=0, ymax=3.3,
      axis lines=left, ytick=\empty, xtick={0,0.5,1}, font=\scriptsize,
      xlabel={$s(x)$ trên validation}, xlabel style={yshift=3pt}, every axis plot/.style={thick, samples=80}]
    \addplot[domain=0:1, okgreen!80!black, fill=okgreen!12] {2.7*exp(-((x-0.18)^2)/(2*0.10^2))} \closedcycle;
    \addplot[domain=0:1, badred, fill=badred!10, opacity=0.9] {1.3*exp(-((x-0.62)^2)/(2*0.17^2))} \closedcycle;
    \draw[warnorange, very thick, dashed] (axis cs:0.31,0) -- (axis cs:0.31,3.2);
    \node[font=\scriptsize, text=warnorange, anchor=west] at (axis cs:0.32,3.0) {$t_{90}$};
    \node[font=\tiny, text=okgreen!70!black] at (axis cs:0.15,2.95) {ảnh âm};
    \node[font=\tiny, text=badred] at (axis cs:0.72,1.55) {ảnh dương};
  \end{axis}
  \end{scope}
\end{tikzpicture}
\caption{Cách tính điểm liên tục $s(x)$ cho VLM và chọn ngưỡng. Trái: hai câu trả lời ép buộc dùng chung ảnh,
câu hỏi và tiền tố; hiệu log-xác suất qua hàm sigmoid cho $s(x)$. Phải (minh hoạ, không phải kết quả): $t_{90}$
đặt sao cho 90\% ảnh âm của validation có điểm dưới ngưỡng; ảnh dương có điểm trên ngưỡng được tính là phát hiện.}
\label{fig:yesno}
\end{figure}

\paragraph{Câu trả lời không parse được.} Khi tính sensitivity và specificity thô, câu trả lời không parse được
tính là ``no''. Riêng các chỉ số chế độ báo bình thường (tỷ lệ ``no'' trên ca dương, ``yes'' trên ca âm) chỉ tính
trên câu trả lời parse được, và tỷ lệ không parse được được báo riêng. Mọi chỉ số được tính lại từ văn bản thô bằng
parser hiện hành.

\paragraph{Sinh văn bản.} C0 dùng giải mã tham lam (nhiệt độ 0) cho dòng kết quả chính, tối đa 256 token. Hai lần
chạy ở nhiệt độ 0,7 chỉ dùng để đo độ ổn định: tỷ lệ đồng thuận trên các ảnh parse được ở mọi lần chạy.

\subsection{Huấn luyện}

\begin{itemize}
  \item SFT với loss chỉ tính trên token câu trả lời; tối ưu AdamW, learning rate $10^{-4}$, lịch cosine,
        warmup 3\%, 3 epoch, batch hiệu dụng 8 (batch 1 $\times$ tích luỹ gradient 8), bf16, gradient checkpointing.
  \item C1, C2 dùng tỷ lệ tự nhiên sau khi giảm bớt ảnh âm (tối đa 4 ảnh âm cho mỗi ảnh dương, cờ
        \code{--neg-per-pos 4}); C3 thêm oversampling ca dương 1:1 (\code{--balance}) và trọng số loss $w_{+} = 2$
        cho ca dương (\code{--pos-weight 2.0}, cố định theo \code{RUNBOOK.md}).
        $w_{+}$ áp lên mọi token đáp án của ca dương, kể cả token box. Loss được chuẩn hoá theo tổng số token nhãn
        trong cả cửa sổ tích luỹ gradient, để $w_{+}$ vẫn có tác dụng khi batch bằng 1 (Hình~\ref{fig:sft}).
  \item Mỗi cấu hình C1--C3 chạy 3 seed; C4 chạy 1--3 seed; B1, B2 mỗi loại 3 seed.
  \item Ngưỡng chọn trên validation; mỗi tập test chỉ dùng một lần. Cấu hình áp cho C4 và cho
        kiểm toán E5 được chọn theo \code{val\_sens90} (sensitivity trên validation tại ngưỡng $t_{90}$ của chính
        lần chạy đó) trung bình qua các seed, đọc từ \code{val\_selection.csv}; không bao giờ dựa trên test.
\end{itemize}

Hình~\ref{fig:sft} cho thấy một mẫu huấn luyện sau khi ghép chuỗi. Toàn bộ token của mẫu chat (token ảnh, câu
hỏi, tiền tố vai trò) được gán nhãn $-100$ nên không đóng góp vào loss. Chỉ token của câu trả lời JSON được tính.
Với một cửa sổ tích luỹ gradient $\mathcal{W}$ gồm 8 mẫu, loss là
\begin{equation}
  \mathcal{L}_{\mathcal{W}} = \frac{\sum_{i \in \mathcal{W}} w_i \sum_{t \in \mathcal{A}_i} -\log p_\theta(y_{i,t} \mid y_{i,<t}, x_i, q)}
  {\sum_{i \in \mathcal{W}} |\mathcal{A}_i|},
  \qquad w_i = \begin{cases} w_{+} & \text{nếu ca dương (C3)} \\ 1 & \text{ngược lại,} \end{cases}
\end{equation}
trong đó $\mathcal{A}_i$ là tập vị trí token đáp án của mẫu $i$. Ở C1 và C2, mọi $w_i = 1$.

\begin{figure}[H]
\centering
\begin{tikzpicture}[x=1cm,y=1cm,
    tok/.style={draw=black!45, minimum height=0.62cm, inner sep=1.5pt, font=\scriptsize\ttfamily, anchor=west},
    msk/.style={tok, fill=black!10, text=black!55},
    ans/.style={tok, fill=uitblue!15, draw=uitblue},
    ansp/.style={tok, fill=badred!12, draw=badred}]
  % ---- ca dương ----
  \node[font=\footnotesize\bfseries, anchor=east] at (-0.1,0) {Ca dương};
  \node[msk, minimum width=1.45cm] (a1) at (0,0) {<image> \dots};
  \node[msk, minimum width=2.75cm] (a2) at (a1.east) {Is there a lung nodule \dots};
  \node[msk, minimum width=1.05cm] (a3) at (a2.east) {model:};
  \node[ansp] (a4) at (a3.east) {\{\dq{}nodule\dq{}: \dq{}};
  \node[ansp] (a5) at (a4.east) {yes};
  \node[ansp] (a6) at (a5.east) {\dq{}, \dq{}boxes\dq{}: [[412, 230, 468, 281]]\}};
  % ---- ca âm ----
  \node[font=\footnotesize\bfseries, anchor=east] at (-0.1,-1.25) {Ca âm};
  \node[msk, minimum width=1.45cm] (b1) at (0,-1.25) {<image> \dots};
  \node[msk, minimum width=2.75cm] (b2) at (b1.east) {Is there a lung nodule \dots};
  \node[msk, minimum width=1.05cm] (b3) at (b2.east) {model:};
  \node[ans] (b4) at (b3.east) {\{\dq{}nodule\dq{}: \dq{}};
  \node[ans] (b5) at (b4.east) {no};
  \node[ans] (b6) at (b5.east) {\dq{}, \dq{}boxes\dq{}: []\}};
  % ---- chú thích ----
  \draw[decorate, decoration={brace, amplitude=4pt, mirror}] (0,-1.7) -- (a3.east |- 0,-1.7)
     node[midway, below=5pt, fnote] {prompt: nhãn $=-100$, không tính loss};
  \draw[decorate, decoration={brace, amplitude=4pt}] (a4.west |- 0,0.42) -- (a6.east |- 0,0.42)
     node[midway, above=5pt, fnote, text=badred] {token đáp án, trọng số $w_{+}=2$ ở C3 (kể cả token box)};
  \draw[decorate, decoration={brace, amplitude=4pt, mirror}] (b4.west |- 0,-1.7) -- (b6.east |- 0,-1.7)
     node[midway, below=5pt, fnote, text=uitblue] {token đáp án, trọng số 1};
\end{tikzpicture}
\caption{Định dạng một mẫu huấn luyện và mặt nạ loss (minh hoạ; số token ảnh và mẫu chat phụ thuộc mô hình).
Ô xám: token ảnh, câu hỏi và tiền tố vai trò, bị loại khỏi loss. Ô màu: token đáp án được tính loss; ở C3, mọi
token đáp án của ca dương nhân với $w_{+}$. Loss chia cho tổng số token nhãn của cả cửa sổ tích luỹ gradient.}
\label{fig:sft}
\end{figure}

\paragraph{Ngân sách tính toán.} Huấn luyện trên Colab. MedGemma chỉ chạy được trên GPU có compute capability
$\ge 8$ (L4 hoặc A100, bf16 thật), vì fp16 bị tràn số. Lượt LoRA đầu tiên dùng để đo thời gian. Nếu một lượt
MedGemma-4B vượt 3 giờ, cắt giảm theo thứ tự: giảm số ảnh âm (giữ toàn bộ ảnh dương), giảm độ phân giải (chỉ với
Qwen2.5-VL; processor của MedGemma cố định 896\,px), giảm còn 2 epoch, giảm C1 còn 1 seed.

\subsection{Kiểm toán ``nhìn đúng chỗ''}
\label{sec:audit}

Một trợ lý trả lời đúng nhưng nhìn sai chỗ (ví dụ dựa vào dấu hiệu máy chụp hay nhãn chữ trên phim) là không an
toàn~\cite{bernstein2023incorrect,lee2023humanai}. Kiểm toán chạy trên mọi ảnh dương của \code{test} (VinBigData)
và \code{test\_jsrt}, với ba phép thử ở Bảng~\ref{tab:audit} và Hình~\ref{fig:occlusion}. Một ảnh được coi là
``yes'' khi $s(x) \ge t_{90}$, với $t_{90}$ lấy từ \code{thresholds.json} của đúng cấu hình đó (seed 0, run 0).

\begin{table}[H]
\centering
\caption{Các phép thử phụ thuộc thị giác.}
\label{tab:audit}
\small
\begin{tabularx}{\textwidth}{@{}l L L@{}}
\toprule
\textbf{Mã} & \textbf{Phép thử} & \textbf{Hành vi mong đợi} \\
\midrule
T1 & Che nốt: box nới 10\% (5\% mỗi cạnh), tô bằng giá trị trung bình của ảnh & Đáp án đổi từ ``yes'' sang ``no'' \\
T2 & Che vùng ngẫu nhiên cùng kích thước với box đã nới, không giao vùng T1, $\ge 90\%$ diện tích nằm trong mặt nạ
phổi (PSPNet của TorchXRayVision, dự phòng bằng mặt nạ heuristic) & Đáp án giữ nguyên \\
T3 & Thay bằng một ảnh âm chọn ngẫu nhiên trong cùng tập test & ``no'' (báo tỷ lệ ``yes'') \\
\bottomrule
\end{tabularx}
\end{table}

\begin{figure}[H]
\centering
\begin{tikzpicture}[x=1cm,y=1cm]
  \newcommand{\cxr}[2]{% ảnh CXR giản lược tại (#1,#2), cỡ 3x3.4
    \fill[black!82, rounded corners=4pt] (#1,#2) rectangle ++(3,3.4);
    \fill[black!45] (#1+0.35,#2+0.55) .. controls (#1+0.3,#2+2.6) and (#1+0.9,#2+3.1) .. (#1+1.35,#2+3.0)
                    -- (#1+1.35,#2+0.75) -- cycle;
    \fill[black!45] (#1+2.65,#2+0.55) .. controls (#1+2.7,#2+2.6) and (#1+2.1,#2+3.1) .. (#1+1.65,#2+3.0)
                    -- (#1+1.65,#2+0.75) -- cycle;}
  % ---- ảnh gốc ----
  \cxr{0}{0}
  \fill[white!85!black] (0.78,2.05) circle (0.17);
  \draw[okgreen, thick] (0.55,1.82) rectangle (1.01,2.28);
  \node[fnote, text=black, font=\footnotesize\bfseries] at (1.5,-0.3) {Ảnh gốc};
  \node[fnote] at (1.5,-0.65) {$s \ge t_{90}$ (``yes'')};
  % ---- T1 ----
  \cxr{3.8}{0}
  \fill[black!62] (4.33,1.80) rectangle (4.83,2.30);
  \draw[badred, thick, dashed] (4.33,1.80) rectangle (4.83,2.30);
  \node[fnote, text=black, font=\footnotesize\bfseries] at (5.3,-0.3) {T1: che nốt};
  \node[fnote] at (5.3,-0.65) {box nới 10\%, tô giá trị TB};
  % ---- T2 ----
  \cxr{7.6}{0}
  \fill[white!85!black] (8.38,2.05) circle (0.17);
  \draw[okgreen, thick] (8.15,1.82) rectangle (8.61,2.28);
  \draw[badred, thin, dashed] (8.13,1.80) rectangle (8.63,2.30);
  \fill[black!62] (9.75,1.15) rectangle (10.25,1.65);
  \draw[warnorange, thick, dashed] (9.75,1.15) rectangle (10.25,1.65);
  \node[fnote, text=black, font=\footnotesize\bfseries] at (9.1,-0.3) {T2: che ngẫu nhiên};
  \node[fnote] at (9.1,-0.65) {cùng cỡ, trong phổi, không giao T1};
  % ---- T3 ----
  \cxr{11.4}{0}
  \node[fnote, text=black, font=\footnotesize\bfseries] at (12.9,-0.3) {T3: tráo ảnh};
  \node[fnote] at (12.9,-0.65) {ảnh âm khác trong tập};
  % ---- công thức ----
  \node[draw=uitblue, fill=softgray, rounded corners=2pt, align=center, font=\small, text width=14.0cm, inner sep=5pt]
     at (7.2,-2.35) {$\displaystyle
     \Delta_{\text{occ}} = \underbrace{\frac{\#\{\text{yes}\rightarrow\text{no ở T1}\}}{\#\{\text{yes ở ảnh gốc}\}}}_{\text{flip\_t1}}
     - \underbrace{\frac{\#\{\text{yes}\rightarrow\text{no ở T2}\}}{\#\{\text{yes ở ảnh gốc}\}}}_{\text{flip\_t2}}$
     \quad trên cùng $n_{\text{T2}}$ ảnh dương lấy mẫu được vùng T2};
\end{tikzpicture}
\caption{Minh hoạ ba phép thử của kiểm toán ``nhìn đúng chỗ'' (ảnh giản lược). Khung xanh: box nốt tham chiếu;
khung đỏ nét đứt: vùng T1 (box nới 10\%); khung cam nét đứt: vùng T2 cùng kích thước, $\ge 90\%$ nằm trong mặt nạ
phổi. $\Delta_{\text{occ}}$ là hiệu tỷ lệ đổi đáp án giữa T1 và T2, tính bắt cặp trên cùng tập ảnh.}
\label{fig:occlusion}
\end{figure}

Chỉ số chính là $\Delta_{\text{occ}} = \text{tỷ lệ đổi đáp án ở T1} - \text{tỷ lệ đổi ở T2}$. Tỷ lệ đổi đáp án là
phần các ảnh có đáp án ``yes'' ở ảnh gốc chuyển thành ``no'' sau khi che. Cả hai vế tính trên cùng tập con bắt cặp
gồm các ảnh lấy mẫu được vùng T2, nên phải báo kèm số ảnh $n_{\text{T2}}$ (khoá \code{n\_t2\_ok}, \code{delta} trong
file \code{audit/*.json}). Tỷ lệ đổi ở T1 trên toàn bộ ca dương (\code{flip\_t1\_all}) chỉ để tham khảo.
$\Delta_{\text{occ}}$ lớn cho thấy quyết định dựa vào vùng nốt; $\Delta_{\text{occ}} \approx 0$ cho thấy mô hình
nhạy với mọi nhiễu loạn chứ không riêng nốt. T3 báo tỷ lệ ``yes'' trên ảnh tráo (\code{yes\_rate\_t3}), mong đợi
gần $1 - 0{,}90 = 0{,}10$. Kết quả được phân tầng theo kích thước nốt và theo độ mờ JSRT 1--5. Với VinBigData và NIH,
kích thước chia theo tertile diện tích box tương đối (ngưỡng lấy từ ca dương của tập train VinBigData, lưu ở
\code{e0\_data\_edges.json}); với JSRT, chia theo mm: $\le 10$, 10--20 và $> 20$\,mm. Kiểm toán áp cho
C0-MedGemma, cấu hình LoRA tốt nhất theo \code{val\_selection.csv} (adapter seed 0) và B2 (trọng số seed 0).
